\section{The notebooks' favourite formula}
\label{sec:mezzo}

The notebook \texttt{mate.dfw} comes back some forty times to the same
identity,
\begin{equation}\label{eq:mezzo}
  \frac{\cos^2x-\cos^2y}{\cos2x-\cos2y}=\frac12 ,
\end{equation}
true for every $x$ and $y$ with $\cos2x\neq\cos2y$, because
$\cos2\theta=2\cos^2\theta-1$ and hence $\cos2x-\cos2y=2(\cos^2x-\cos^2y)$
(\riga{mate.dfw}{2}). The interesting part is the way the notebook puts it to
the test.

\subsection{The variants}

\begin{itemize}[leftmargin=1.4em,itemsep=2pt]
  \item With sines the ratio is $-\tfrac12$, because
        $\sin^2\theta=(1-\cos2\theta)/2$ changes the sign (\riga{mate.dfw}{4});
        with hyperbolic cosines it is $\tfrac12$ again, because
        $\cosh2\theta=2\cosh^2\theta-1$ has the same form (\riga{mate.dfw}{3}).
  \item The identity holds for arbitrary complex arguments: lines
        \righe{mate.dfw}{19}{31} test it with $x\pm\ii y$, with $x\pm ky$, with
        $3x+ky$ and $x$, and line \#6 rewrites it with complex exponentials.
  \item With the triple-angle formula $\cos3\theta=4\cos^3\theta-3\cos\theta$
        the constant halves again (\riga{mate.dfw}{5}):
        \begin{equation}\label{eq:quarto}
          \frac{\cos^3x\,\cos y-\cos^3y\,\cos x}{\cos3x\,\cos y-\cos3y\,\cos x}=\frac14 .
        \end{equation}
        Multiplying by the other cosine removes the term $-3\cos\theta$.
  \item With two weights, the ratio
        $(\cos^2x-k\cos^2y)/(\cos2x-l\cos2y)$ equals $\tfrac12$ for every $x$ if
        and only if $2k\cos^2y=1+l\cos2y$, that is $\cos2y=(1-k)/(k-l)$: this is
        the relation Derive finds in \righe{mate.dfw}{53}{60}.
  \item Finally, \eqref{eq:mezzo} becomes a ``machine for $\pi$''
        (\righe{mate.dfw}{67}{71}): since $(1/2)!=\Gamma(3/2)=\sqrt\pi/2$,
        \[
          4\left[\left(\frac{\cos^2x-\cos^2y}{\cos2x-\cos2y}\right)!\,\right]^2=\pi
        \]
        for every pair of angles.
\end{itemize}

\subsection{The general form}

The constants $\tfrac12$ and $\tfrac14$ are the reciprocals of the leading
coefficients of the Chebyshev polynomials $T_2(c)=2c^2-1$ and $T_3(c)=4c^3-3c$,
with $\cos n\theta=T_n(\cos\theta)$. With only two angles the trick stops at
degree three: $T_4(c)=8c^4-8c^2+1$ has two lower-degree terms to remove, and a
combination of two cosines is not enough. With more angles, however, the rule
holds in every degree.

\begin{prop}\label{prop:chebyshev}
Let $n\ge2$, $m=\lfloor n/2\rfloor+1$, let $x_1,\dots,x_m$ be angles and
$c_i=\cos x_i$. Then
\[
  \det\bigl[\cos nx_i,\ c_i^{\,n-2},\ c_i^{\,n-4},\ \dots,\ c_i^{\,n-2m+2}\bigr]
  =2^{n-1}\det\bigl[c_i^{\,n},\ c_i^{\,n-2},\ c_i^{\,n-4},\ \dots,\ c_i^{\,n-2m+2}\bigr],
\]
where each expression denotes the $m\times m$ matrix whose rows are indexed by
$i$. The determinant on the right is nonzero if the $c_i^2$ are distinct and,
for odd $n$, nonzero.
\end{prop}

\begin{proof}
The polynomial $T_n$ has the parity of $n$:
$T_n(c)=2^{n-1}c^n+\sum_{j\ge1}a_jc^{n-2j}$, and the powers $c^{n-2j}$ with
$j\ge1$ and $n-2j\ge0$ are exactly the other $m-1$ columns. Subtracting them
from the first column with coefficients $a_j$ leaves the determinant unchanged,
and what remains is $2^{n-1}c_i^n$. On the right, for even $n$ the columns are
the powers $c_i^{2k}$ with $0\le k<m$, and the determinant is a Vandermonde in
the $c_i^2$; for odd $n$ one factors $c_i$ out of each row and finds the same
Vandermonde.
\end{proof}

For $n=2$ and $n=3$ there are two angles, and the proposition gives back
exactly \eqref{eq:mezzo} and \eqref{eq:quarto}. For $n=4$ three angles are
needed:
\[
  \det\begin{bmatrix}
    \cos4x_1 & c_1^2 & 1\\
    \cos4x_2 & c_2^2 & 1\\
    \cos4x_3 & c_3^2 & 1
  \end{bmatrix}
  =8\det\begin{bmatrix}
    c_1^4 & c_1^2 & 1\\
    c_2^4 & c_2^2 & 1\\
    c_3^4 & c_3^2 & 1
  \end{bmatrix}.
\]

Identity \eqref{eq:mezzo} is also the source of the trap described in
Section~\ref{sec:refusi}: since it holds for every value of a parameter, the
numerical solver asked to find that value still answers with a number.

\paragraph{What remains open.}
Nothing: the proposition is an elementary observation about Chebyshev
polynomials. It is, however, the general form of the rule that the notebook
glimpses in degrees two and three.
